\section{A finite-dimensional categorical geometry primer}
\label{sec:book-categorical-geometry-primer}

This primer collects the notation needed to reread the preceding table calculation and follow the
exact construction below.  All spaces in this chapter are finite-dimensional.  The notation $L^2(p)$ therefore means an
ordinary vector space equipped with weights, rather than an infinite-dimensional function space.
For a full-support categorical reference $p=(p(1),\ldots,p(q))^{\mathsf T}$, write
\begin{equation}
D=\operatorname{diag}(p),
\qquad
\langle f,g\rangle_p
=\sum_{a=1}^q p(a)f(a)g(a)
=f^{\mathsf T}Dg.
\label{eq:book-primer-weighted-inner-product}
\end{equation}
A vector $f\in\mathbb R^q$ is a numerical score on the $q$ residue categories.  Its reference
mean is $p^{\mathsf T}f$, and the centered contrast space is
\begin{equation}
\mathcal H^0
=\{f\in\mathbb R^q:p^{\mathsf T}f=0\}.
\label{eq:book-primer-centered-subspace}
\end{equation}
It has dimension $q-1$: adding the same number to every category supplies no contrast.

The matrix
\begin{equation}
P=I-\mathbf1p^{\mathsf T}
\label{eq:book-primer-weighted-projection}
\end{equation}
subtracts the weighted mean, since $Pf=f-\mathbf1(p^{\mathsf T}f)$.  It obeys
$P^2=P$ and $P^{\mathsf T}D=DP$, so it is the orthogonal projector onto $\mathcal H^0$ in the
$p$-weighted inner product.  For a uniform reference this is ordinary mean subtraction.  For a
nonuniform reference, orthogonality is measured with the declared category frequencies.

Whitening converts the weighted geometry into Euclidean geometry.  The coordinate change
\begin{equation}
f\longmapsto \widetilde f=D^{1/2}f
\end{equation}
satisfies $\|f\|_p=\|\widetilde f\|_2$, and the projector becomes
\begin{equation}
Q=D^{1/2}PD^{-1/2}
=I-\sqrt{p}\sqrt{p}^{\,\mathsf T}.
\label{eq:book-primer-whitened-projection}
\end{equation}
For endpoint references $p_i,p_j$, a centered pair block $G$ has weighted norm
\begin{equation}
\|G\|_{p_i,p_j}^2
=\sum_{a,b}p_i(a)p_j(b)G(a,b)^2
=\|D_i^{1/2}GD_j^{1/2}\|_F^2.
\label{eq:book-primer-weighted-block-norm}
\end{equation}

The singular-value decomposition of the whitened block is
\begin{equation}
\widetilde G
=D_i^{1/2}GD_j^{1/2}
=\sum_{r=1}^{q-1}\sigma_r u_rv_r^{\mathsf T}.
\label{eq:book-primer-block-svd}
\end{equation}
The same block induces the edge map $\mathcal Kf=GD_jf$ from source contrasts to target contrasts;
in whitened coordinates its matrix is $\widetilde G$.  Each $v_r$ is a whitened source-category
contrast, $u_r$ is the whitened target contrast it produces, and $\sigma_r$ is the gain along that
pair of directions.  The Hilbert--Schmidt norm keeps their total squared strength,
\begin{equation}
\|\mathcal K\|_{\mathrm{HS}}^2
=\|\widetilde G\|_F^2
=\sum_r\sigma_r^2,
\label{eq:book-primer-hs-norm}
\end{equation}
while discarding the individual source and target directions.

Finally, the quotient language records a simple equivalence relation.  Two raw tables represent
the same pair interaction when
\begin{equation}
T'-T=u\mathbf1^{\mathsf T}+\mathbf1v^{\mathsf T}
+c\mathbf1\mathbf1^{\mathsf T}.
\label{eq:book-primer-table-equivalence}
\end{equation}
They differ only by target-only, source-only, and constant terms, and double projection returns the
same centered block.  A quotient class is the collection of all raw tables related in this way;
the reference $p_i,p_j$ selects one centered representative of that class.

Direct sums assemble one contrast vector per node, while tensor products later encode contrasts
involving several source coordinates at once.
